\section{绪论}

\subsection{研究背景与意义}

随着物联网、工业互联网、可穿戴医疗与智慧环境监测系统的快速发展，终端节点对“长期自治供电 + 高可信感知”的要求不断提高。对于这类系统而言，前端传感与数据转换链路往往是总能耗与系统精度的关键瓶颈。一方面，节点常由纽扣电池、薄膜电池或能量采集模块供电，平均功耗预算通常处于微瓦到毫瓦量级；另一方面，环境参数、生理信号和工况状态的微小变化又要求读出电路具备足够分辨率、线性度和温漂稳定性。由此，低功耗高精度接口电路设计成为集成电路研究中的重要方向\cite{tang2020zoomcdc,li2021humidityarray,li2024crcls}.

在传统模拟前端中，静态偏置 OTA 长期承担积分、缓冲和放大任务。这一路径在高性能场景下成熟可靠，但在低电压与超低功耗条件下面临明显挑战：(1) 为获得高增益和高线性，通常需要较大偏置电流，导致待机与运行能耗偏高；(2) 先进工艺下晶体管本征增益下降，静态 OTA 维持高环路增益的代价进一步上升；(3) 供电电压降低后，级联结构与信号摆幅受限，性能提升空间收窄。因而，“以持续静态电流换取精度”的传统模式，越来越难同时满足边缘传感系统对能效、面积和鲁棒性的综合要求。

动态放大器（Dynamic Amplifier）由此受到广泛关注。其核心思想是利用离散时间相位中的瞬态过程实现信号增益，而非全时偏置放大。对于低带宽传感读出、离散时间积分和噪声整形数据转换场景，这一机制天然契合“按需计算”的系统节奏。以近年的代表性工作为例：在低频高精度 ADC 中，基于 floating-inverter 的动态放大路径可在微瓦级功耗下获得较高 SNDR\cite{wang2021sefiamod,meng2023jssc_mash}；在 CDC 与湿度传感接口中，动态放大器阵列与 zoom 架构协同可显著改善能效指标\cite{li2021humidityarray,tang2020zoomcdc}。这些结果说明，动态放大器不仅是“替代 OTA 的局部电路技巧”，更是面向低功耗系统架构优化的关键支点。

尽管如此，动态放大器并非没有代价，其工程实现存在一组典型矛盾。其优势主要体现在：(1) 静态电流近零，适合间歇采样与低占空比工作；(2) 与开关电容和数字时序天然兼容，便于系统级时钟编排；(3) 在特定带宽范围内能实现较优能效。与此同时，其不足也较为突出：(1) 增益形成依赖时钟边沿和充放电轨迹，对相位偏差、抖动与非重叠时序敏感；(2) 动态节点较多，开关注入、寄生耦合与电容失配会直接映射为非线性与噪声；(3) 在低电压下，输入共模和输出摆幅进一步受限，若缺乏配套补偿机制，易出现闭环增益误差放大。换言之，动态放大器的价值不在于“单点提升某一指标”，而在于通过结构、时序与系统算法的联合设计，把功耗-线性-鲁棒性折中推向更优边界\cite{tang2020ns_sar,tang2020zoomcdc}.

围绕上述问题，近年来“相关电平移位（Correlated Level Shifting, CLS）+ 动态放大器”成为重要研究方向。CLS 的核心是在相关时钟相位内执行受控电平重定位，使部分误差项在后续处理中具备可抑制性，从而改善低压条件下的有效增益与线性。已有工作表明，CLS 可在离散时间 MASH ADC、inverter-based OTA 以及传感 CDC 前端中带来可观收益\cite{hu2022mashcls,meng2023jssc_mash,liu2020pingpongcls,li2024crcls}。其中，面向湿度传感的 CR-CLS（Charge-Redistribution CLS）进一步展示了“电荷重分配 + 相关电平移位 + FIA + Zoom CDC”的跨层协同潜力，为低功耗高精度传感读出提供了新的实现路径\cite{li2024crcls}.

从应用需求看，未来大量边缘节点并不追求最高采样速率，而更强调长期稳定、低能耗与可批量制造，这恰好与动态放大器及其改进技术的演进方向一致。对本课题而言，研究基于 CLS 与低压浮动反相放大器（LVFIA）的高能效动态放大电路，既具有明确的理论价值，也具有直接工程意义：在理论层面，可深化对低压动态放大误差机理与相关抑制机制的理解；在工程层面，可为低功耗传感接口与高分辨离散时间转换器提供可复用的设计方法，支撑高性能边缘感知芯片的实现。
